% !TEX root = ../main.tex
\chapter*{Introduction générale}
\addcontentsline{toc}{chapter}{Introduction générale}
\markboth{Introduction générale}{Introduction générale}
\justifying

La gestion des mobilités internationales dans les établissements d'enseignement supérieur implique la manipulation de nombreuses informations : les étudiants, les universités partenaires, les accords de mobilité, les candidatures, les critères d'éligibilité ou encore les quotas disponibles. Ces informations sont souvent réparties entre plusieurs outils et systèmes, ce qui peut rendre leur gestion plus complexe pour les services chargés des relations internationales.

Dans ce contexte, certaines tâches sont encore réalisées manuellement, notamment la collecte et la consolidation des données provenant de différentes sources, le suivi des candidatures et des mobilités, ainsi que la gestion des affectations des étudiants. Cette organisation peut entraîner une perte de temps, des difficultés de suivi et un manque de visibilité, aussi bien pour les étudiants que pour les gestionnaires.

C'est dans ce contexte que s'inscrit ce projet de fin d'études, réalisé dans le cadre du cycle ingénieur en génie informatique de l'École Nationale d'Ingénieurs de Carthage, au sein du service des Relations Internationales de Toulouse INP-ENSEEIHT. L'objectif est de concevoir et de développer une plateforme web permettant de centraliser et de faciliter la gestion des mobilités internationales des étudiants. La plateforme prend notamment en charge les mobilités académiques sortantes, les stages internationaux, les mobilités complémentaires et l'accueil d'étudiants étrangers, tout en s'appuyant sur les données provenant des systèmes institutionnels existants, tels que MoveON, Pégase et Eudonet. Elle vise également à automatiser certaines étapes du processus, notamment la gestion des candidatures et l'affectation des étudiants.

Ce rapport détaille le travail réalisé dans le cadre de ce projet. Il s'ouvre sur cette introduction générale, puis se compose des chapitres suivants :
\begin{itemize}[label=•]
    \item Le chapitre~\ref{cha:contexte}, \textit{Contexte, problématique et positionnement du projet}, présente l'organisme d'accueil, la manière dont les mobilités internationales sont gérées aujourd'hui, la problématique qui en découle et le positionnement du projet par rapport aux solutions existantes.
    \item Le chapitre~\ref{cha:besoins-technologies}, \textit{Besoins fonctionnels et choix technologiques}, détaille les besoins fonctionnels du système ainsi que les choix technologiques retenus pour y répondre.
    \item Le chapitre~\ref{cha:conception}, \textit{Conception du système}, présente la conception générale du système : son architecture, son modèle de données et les flux qui le traversent.
    \item Le chapitre~\ref{cha:algo-fsm}, \textit{Choix algorithmiques et machines à états}, justifie les choix algorithmiques structurants, l'algorithme d'affectation et la méthode de répartition des quotas, ainsi que les machines à états qui sécurisent le cycle de vie des entités du système.
    \item Le chapitre~\ref{cha:realisation}, \textit{Réalisation}, présente la mise en œuvre du projet à travers ses six sprints, menés selon la méthodologie Scrum. Chaque sprint y est détaillé par son backlog, sa conception, son implémentation et les tests qui la valident.
    \item Le chapitre~\ref{cha:tests-validation}, \textit{Tests et validation du système}, revient sur la stratégie de tests et de validation retenue à l'échelle du projet.
\end{itemize}

Le rapport se clôture par une conclusion générale qui résume les principaux résultats obtenus et présente les perspectives d'évolution du système.
